% Part: first-order-logic
% Chapter: tableaux
% Section: rules-and-proofs

\documentclass[../../../include/open-logic-section]{subfiles}

\begin{document}

\iftag{FOL}
      {\olfileid{fol}{tab}{rul}}
      {\olfileid{pl}{tab}{rul}}

\olsection{Regels en \usetoken{P}{tableau}}

!!^a{tableau} geeft een systematisch overzicht van de mogelijke
manieren waarop !!a{sentence} in !!a{structure} waar of onwaar kan
zijn. De bouwstenen van een tableau zijn !!{signed formula}s:
!!{sentence}s met een waarheidswaarde als ``teken'', hetzij $\True$,
hetzij~$\False$. Deze !!{formula}s zijn gerangschikt in een (naar
beneden groeiende) boom.

\begin{defn}
  Een \emph{!!{signed formula}} is een paar dat bestaat uit een
  waarheidswaarde en !!a{sentence}, dus een van de volgende twee
  vormen:
  \[
  \sFmla{\True}{!A} \text{ of } \sFmla{\False}{!A}.
  \]
\end{defn}

Intuïtief kunnen we $\sFmla{\True}{!A}$ lezen als ``$!A$ zou waar
kunnen zijn'' en $\sFmla{\False}{!A}$ als ``$!A$ zou onwaar kunnen
zijn'' (in een zekere !!{structure}).

Elke !!{signed formula} in de boom is ofwel een \emph{aanname} (de
aannamen staan helemaal bovenaan in de boom), ofwel verkregen uit
!!a{signed formula} erboven door middel van een van de
afleidingsregels. Voor elke mogelijke !!{main operator} van de
voorafgaande !!{formula} zijn er twee regels: een voor het geval dat
het teken~$\True$ is en een voor het geval dat het teken~$\False$ is.
Sommige regels laten de boom vertakken; andere voegen alleen
!!{signed formula}s aan de tak toe. Een regel hoeft niet op de
onmiddellijk voorafgaande !!{signed formula} te worden toegepast. Hij
kan (en moet vaak) worden toegepast op een willekeurige !!{signed
formula} in de tak van de wortel tot de plaats waar de regel wordt
toegepast.

Een tak is \emph{gesloten} als hij zowel $\sFmla{\True}{!A}$ als
$\sFmla{\False}{!A}$ bevat. !!^a{tableau} is gesloten als elke tak
gesloten is. Volgens de intuïtieve interpretatie beschrijft elke tak
een gezamenlijke mogelijkheid, maar $\sFmla{\True}{!A}$ en
$\sFmla{\False}{!A}$ zijn niet gezamenlijk mogelijk. Als een tak
gesloten is, is de mogelijkheid die hij beschrijft met andere woorden
uitgesloten. Dit betekent in het bijzonder dat een gesloten
!!{tableau} alle mogelijkheden uitsluit om tegelijk elke aanname van
de vorm $\sFmla{\True}{!A}$ waar en elke aanname van de
vorm~$\sFmla{\False}{!A}$ onwaar te maken.

Een gesloten !!{tableau} \emph{voor $!A$} is een gesloten
!!{tableau} met wortel~$\sFmla{\False}{!A}$. Als zo'n gesloten
!!{tableau} bestaat, zijn alle mogelijkheden waarin~$!A$ onwaar is
uitgesloten; $!A$ moet dus waar zijn in elke !!{structure}.

\end{document}
